\subsection{An inverse low-degree restriction and the torsion it can move to degree three}
\label{arch136:low-generated-inverse}

\paragraph{ARCH-136-LOW-GENERATED: scope and statement.}
Let \(A\) be the free source constructed in
ARCH-136-DEGREE-TWO-FREE, using only its stated horizontal unit-disk
attachments and reduced-representable Tate pairs in degrees
\(q\geq3\). The assertions below also hold if finitely many of
those permitted attachments are performed before the alternating
construction begins. Define the simplicial differential PD subalgebra
\(D(A)\subset A\) generated by all of \(A^0\) and \(A^1\).
Define a second graded subobject by
\[
 \begin{split}
 (A^\dagger)^0&=A^0,\qquad (A^\dagger)^1=A^1,\\
 (A^\dagger)^2&=A^0\,d(A^1)+A^1A^1=D(A)^2,\\
 (A^\dagger)^q&=A^q\quad(q\geq3).
 \end{split}
 \tag{LG.1}
\]
All products and sums in (LG.1) mean their submodules inside \(A\).
Then \(A^\dagger\) is an actual simplicial strict differential
PD subalgebra, with its original maps to all targets, and in every
weight and simplex
\[
 H^0(A^\dagger)=S,\qquad H^1(A^\dagger)=H^2(A^\dagger)=0.
 \tag{LG.2}
\]
It has the same horizontal rows as \(A\) except in form degree two,
and \(H^q(A^\dagger)=H^q(A)\) for \(q\geq4\). Its changed
degree-three group fits into
\[
 0\longrightarrow
 \frac{A^2}{D(A)^2+Z^2(A)}
 \longrightarrow H^3(A^\dagger)
 \longrightarrow H^3(A)\longrightarrow0,
 \tag{LG.3}
\]
interpreted in each weight and simplex. Thus all the residue maps
\(\delta_{jk,N}\) vanish on \(A^\dagger\), but this is not yet
an admissible source: the form-two horizontal row and degree-three
vertical cohomology remain to be controlled.

\paragraph{Proof.}
The definition is stable under every simplicial map. It is closed
under the differential and products: the only new product to check
in degree two is a product of two one-forms, which was included.
Multiplication by zero-forms preserves that submodule, and all
products of higher total form degree land in a retained component.
The coefficient and zero-form divided powers are retained. Higher
divided powers of positive even form degree land in degree at least
four, which is also retained. This proves the strict PD-subalgebra
assertion.

We check the stronger vanishing in (LG.2) with the actual free
presentation. Fix a simplex. Choose compatible free complements
at successive horizontal attachments as in
ARCH-136-DEGREE-TWO-FREE. Every new horizontal pair is then a
vertical unit disk. In forms zero and one the resulting generators
consist of zero-form generators \(x\), their independent partners
\(dx\), and further odd form-one generators \(u_i\). Each \(u_i\)
has its own independent even form-two partner \(v_i\), with
\[
                  du_i=N_i v_i+b_i,
                  \qquad N_i\in\mathbb Z\setminus\{0\},
 \tag{LG.4}
\]
where \(b_i\) involves only earlier generators. For horizontal
form-one disks \(N_i=1,b_i=0\); for a degree-three Tate pair this
is its displayed relation, with the sign incorporated in \(b_i\).
Other generators have form degree at least two. Later attachments
do not change the differential or the independence of the older
generators.

Over \(\mathbb Q\), replacing each \(v_i\) by \(du_i\) is a
triangular invertible change of free graded-algebra coordinates.
Every polynomial and its inverse uses only finitely many earlier
generators. Therefore the map from the abstract free differential
PD algebra on the unit disks
\[
       x\longmapsto dx,\qquad u_i\longmapsto y_i,
 \tag{LG.5}
\]
to \(A\), sending \(y_i\) to \(du_i\), is injective after
rationalization. Its integral source is free as an abelian group,
so it is already injective integrally. Its image is exactly
\(D(A)\): it contains every zero- and one-form and their
derivatives, and these generate that image under products and
divided powers. This description does not require the chosen free
complements to be simplicial.

The zero-form unit disks are integrally contractible, and the odd
form-one unit disks have no positive cohomology through degree
three, by the explicit block computation in
ARCH-136-DISK-SUPPORT. Hence \(H^2(D(A))=0\). The components of
\(A^\dagger\) and \(D(A)\) agree through degree two, and the
inclusion of their degree-three components is injective. Their
degree-two cycles and boundaries are consequently identical.
This proves \(H^2(A^\dagger)=0\). The degrees zero and one and
their differential images were unchanged from \(A\), proving
the other two assertions in (LG.2).

The quotient complex \(A/A^\dagger\) is concentrated in form
degree two, with term \(A^2/D(A)^2\). Its cohomology exact
sequence, and \(H^2(A^\dagger)=0\), give (LG.3) and the
asserted identities in degrees at least four. The row statement
follows directly from (LG.1). In particular \(A^\dagger_0=S\).
The residue maps have zero source in every simplex, so their
vanishing follows without any assertion about determinant
operations or a common first skeleton. \(\square\)

\paragraph{ARCH-136-LOW-GENERATED-TORSION: a limitation of arbitrary presentation choices.}
The new group in (LG.3) need not be torsion-free, even in the
weight tail of a completed degree-two-free source. The example
below deliberately uses two avoidable zero-cycle Tate pairs.
It disproves an automatic torsion-freeness inference; it does
not show that a more carefully chosen presentation must fail.

Fix a prime \(p\) and \(m\geq2\). At one simplex consider the
free graded PD algebra over \(S\) with odd generators
\(u,v,t\) of form degree one, even generators \(x,y,w\) of
form degree two, and
\[
 \begin{gathered}
 du=px,\qquad dv=py,\qquad dx=dy=0,\\
 z=xv-uy,\qquad dw=z,\qquad dt=pw-uv.
 \end{gathered}
 \tag{LG.6}
\]
The weights of \(u,v,x,y\) are \(m\), and those of \(t,w\)
are \(2m\). The identities \(dz=0\), \(d(uv)=pz\) verify
\(d^2=0\). The algebra is rationally a tensor product of unit
disks. Its degree-one cohomology is zero, and direct computation
gives \(H^2=(S/pS)x\oplus(S/pS)y\).

Here the degree-three part of \(D(A)\), the algebra generated
in degrees zero and one, is generated over \(S\) by products
of a one-form with one of
\(px,py,pw-uv\), and by \(uvt\). Every coefficient of a
degree-three basis monomial involving \(x,y,w\) is divisible
by \(p\). Thus
\[
 z\notin D(A)^3,
 \qquad pz=d(uv)\in dD(A)^2.
 \tag{LG.7}
\]
The same reasoning shows that the classes
\(T^{[s]}z\), \(s\geq0\), give a copy of
\((S/pS)\langle2m\rangle\) in \(H^3(A^\dagger)\).

We justify that this pattern can survive in an actual final source,
rather than only in this small algebra. Starting from \(S\),
perform two permitted degree-three Tate attachments on
\(\mathbb Z[\Delta[1]]/\mathbb Z[V_1]\), of weight \(m\),
with zero cycle, zero relation cochain, and chosen integer \(p\).
These give the pairs \((u,x)\) and \((v,y)\) in (LG.6) in
simplicial degree one. The cycle \(z\) there has order \(p\),
with relation cochain \(uv\). Attach its degree-three Tate pair
on the same reduced representable, of weight \(2m\), giving
\((t,w)\). All faces of the attaching data in simplex zero
are zero. Send every added generator to zero in the standard
admissible theory \(B\). These are strict permissible maps.
Now perform the full alternating construction of
ARCH-136-DEGREE-TWO-FREE. Its proof applies unchanged from this
finite free initial segment, yielding an actual final source
\(\widehat A\to B\) with all its stated horizontal and higher
vertical tails.

At every subsequent stage, in each fixed simplex, one has
\[
             D(A_{\rm new})\cap A_{\rm old}=D(A_{\rm old}).
 \tag{LG.8}
\]
For a zero-form horizontal disk this follows from its independent
free coordinates. A form-one horizontal disk or a degree-three
Tate pair adds new odd \(u_i\) and independent even \(v_i\)
as in (LG.4). A polynomial in \(u_i,du_i\) with coefficients
in \(D(A_{\rm old})\) can be independent of the new \(u_i,v_i\)
only if all its positive new-coordinate coefficients vanish:
the leading coefficient in a divided power of \(v_i\) is a
nonzero power of \(N_i\), and the old algebra is torsion-free.
Then its constant coefficient belongs to \(D(A_{\rm old})\).
This proves (LG.8) for these attachments. Attachments with both
generators in forms at least two add no generator to \(D(A)\),
so (LG.8) is immediate for them. The statement for arbitrary
simultaneous sets follows by finite monomial support, and the
statement at a union stage follows because every element uses
only finitely many generators.

Consequently none of the old \(T^{[s]}z\) enters
\(D(\widehat A)^3\). It cannot be a boundary from
\((\widehat A^\dagger)^2=D(\widehat A)^2\), while its multiple
by \(p\) remains the boundary \(d(T^{[s]}uv)\). Thus the
displayed copy of \(S/pS\) persists in
\(H^3((\widehat A^\dagger)_1)\) at arbitrarily high weights.
This also checks explicitly that quotienting by \(Z^2\) in
(LG.3) does not remove the example. \(\square\)

\paragraph{The remaining constructive question.}
One may repair the form-two row of \(A^\dagger\) by horizontal
free disks with even lower generator of form degree two. These
extensions are integrally vertically contractible and have maps
to \(A\) using its acyclic form-two row. They can, however,
change higher horizontal rows; no simultaneous bound for all
those rows is supplied here. The example above also shows why
an arbitrary choice of degree-three primitive presentation is
insufficient. A successful use of (LG.1) would need a compatible
choice of actual cohomology generators and their simplicial and
coefficient relations, with controlled degree-three torsion.
The integral common-source question remains open in this work.
